\documentclass[a4paper,11pt]{article}
\usepackage{exam-style-842-unit}
\examinehead{842 信息化概论 · 第 1 单元练习册}{练习二\quad 信息社会与发展战略}
\begin{document}

\unitcover{842 信息化概论 · 单元练习册}{第 1 单元\quad 信息技术与信息社会}{练习二\quad 信息社会与发展战略（共 100 分）}{约 70 分钟}{数据、信息、知识、智慧；信息的分类与特点；信息科学三论；信息技术与信息产业；信息社会与信息化；四次信息技术革命；我国信息化发展战略。}

\begin{knowmap}
\textbf{信息化：}充分利用信息技术、开发利用信息资源、促进信息交流和知识共享、提高经济增长质量、推动经济社会发展转型的历史进程，是动态的、渐进的过程\\[2pt]
\textbf{社会形态：}农业社会（土地·劳动力）$\rightarrow$ 工业社会（资本·机器）$\rightarrow$ 信息社会（战略资源＝信息和知识，主导产业＝知识密集型产业）\\[2pt]
\textbf{信息化特性：}“四化” 智能化·电子化·全球化·非群体化；“四性” 综合性·竞争性·渗透性·开放性\\[2pt]
\textbf{信息社会四特征：}知识型经济 · 网络化社会 · 数字化生活 · 服务性政府\\[2pt]
\textbf{四次信息技术革命：}语言与文字 $\rightarrow$ 造纸与印刷 $\rightarrow$ 电报·电话·电视 $\rightarrow$ 电子计算机与现代通信\\[2pt]
\textbf{认识转变与战略：}重内涵$\rightarrow$重外延；狭义工具论$\rightarrow$计算思维；人机共生$\rightarrow$三元社会；《2006—2020 年国家信息化发展战略》提出 9 项重点
\end{knowmap}

\begin{keyterms}
\textbf{信息化}：充分利用信息技术、开发利用信息资源，推动经济社会发展转型的历史进程。\\
\textbf{信息社会}：脱离工业化社会以后信息起主要作用的社会，战略资源是信息和知识。\\
\textbf{信息社会指数}：国家信息中心提出的信息社会发展水平定量测评指标与体系。\\
\textbf{“四化”}：智能化、电子化、全球化、非群体化。\\
\textbf{“四性”}：综合性、竞争性、渗透性、开放性。\\
\textbf{知识型经济}：人力资源知识化、发展方式可持续化、产业结构软化、经济水平发达。\\
\textbf{网络化社会}：信息基础设施的完备性、社会服务的包容性、社会发展的协调性。\\
\textbf{数字化生活}：生活工具数字化、生活方式数字化、生活内容数字化。\\
\textbf{服务性政府}：科学决策、公开透明、高效治理、互动参与。\\
\textbf{数字化}：将模拟信号通过抽样、量化、编码等环节变换成数字码元。\\
\textbf{计算思维}：面向所有人的普适思维，强调一切皆可计算，不是程序化思维。\\
\textbf{三元世界}：物理世界、信息世界、人类社会组成的动态开放网络社会。\\
\textbf{电子政务}：政府运用信息技术优化管理与服务流程，提供规范、高效的管理和服务。\\
\textbf{数字鸿沟}：不同地区或群体在信息基础设施与信息技术应用水平上的差距。
\end{keyterms}

\begin{contrast}{信息化 / 工业化 / 城镇化（每行依次为：核心内涵·过程特点·相互关系）}
\citem{信息化}{充分利用信息技术，开发利用信息资源，提高经济增长质量}{动态的、渐进的过程，社会经济结构不断变换}{与工业化深度融合，与城镇化、农业现代化同步发展}
\citem{工业化}{机器大生产逐步取代手工劳动，工业在国民经济中占主导地位}{物质生产与资本积累不断扩张的过程}{为信息化提供物质技术基础，与城镇化良性互动}
\citem{城镇化}{人口向城镇聚集，城镇生产生活方式不断扩散}{城乡结构、就业结构持续调整的过程}{与工业化良性互动，与农业现代化相互协调}
\end{contrast}

\begin{contrast}{数字化 / 信息化（每行依次为：核心含义·地位·相互关系）}
\citem{数字化}{将模拟信号经抽样、量化、编码变换成数字码元}{实现信息化的技术手段与必备条件}{为信息化的实现提供有力的技术支撑}
\citem{信息化}{以信息技术为基础、以信息环境为依托，实现信息资源共享和实时交换}{强调信息的决定性作用，把控制和利用信息作为根本要求}{以信息的数字化为必备条件}
\end{contrast}

\begin{contrast}{信息社会 / 工业社会（每行依次为：战略资源·主导产业·社会标志）}
\citem{信息社会}{信息和知识}{知识密集型产业}{以信息技术为基础，以信息价值的生产为中心，以信息产品为标志}
\citem{工业社会}{资本和能源}{制造业等资本密集型产业}{以大机器生产为中心，工业产值占主导}
\end{contrast}

\begin{contrast}{狭义工具论 / 计算思维 / 三元社会模式（每行依次为：核心主张·思维属性·意义或影响）}
\citem{狭义工具论}{把计算机与信息网络仅看作高科技工具}{认知狭隘，视信息科技为辅助性工具学科}{是信息科技向各行各业渗透的最大障碍}
\citem{计算思维}{一切皆可计算，是每个人的基本技能}{概念化思维、人的思维、面向所有人的思维}{促进信息科技在全社会的普及}
\citem{三元社会模式}{物理世界、信息世界、人类社会组成三元世界}{多人、多机、多物组成的动态开放网络社会}{信息科学应研究人机物三者在社会中的信息处理过程}
\end{contrast}

\qtypename{一、单项选择题}{10 题，每题 3 分，共 30 分}
\q 信息社会的战略资源是（\quad）
\choicesline{A. 资本与机器}{B. 信息与知识}{C. 土地与劳动力}{D. 能源与材料}

\q 在信息社会中起主导作用的产业是（\quad）
\choicesline{A. 劳动密集型产业}{B. 资本密集型产业}{C. 知识密集型产业}{D. 资源密集型产业}

\q “充分利用信息技术，开发利用信息资源，促进信息交流和知识共享，提高经济增长质量，推动经济社会发展转型的历史进程”，这一表述指的是（\quad）
\choicesline{A. 数字化}{B. 信息化}{C. 工业化}{D. 城镇化}

\q 1963 年在《信息产业论》中首次提出“信息社会”概念的是日本社会学家（\quad）
\choicesline{A. 梅棹忠夫}{B. 贝尔}{C. 托夫勒}{D. 维纳}

\q 下列不属于信息社会基本特征的是（\quad）
\choicesline{A. 知识型经济}{B. 网络化社会}{C. 数字化生活}{D. 机械化生产}

\q 按信息社会发展指数（ISI）划分，ISI 值在 0.3 以下属于（\quad）
\choicesline{A. 起步期}{B. 转型期}{C. 初级阶段}{D. 高级阶段}

\q 下列关于电子政务的理解，正确的是（\quad）
\choicesline{A. 即政府网上采购}{B. 优化管理与服务}{C. 等同电子商务}{D. 仅限内部办公}

\q 计算思维是（\quad）
\choicesline{A. 程序化思维}{B. 概念化思维}{C. 计算机的思维}{D. 纯数学性思维}

\q 物理世界、信息世界和人类社会组成的三元世界，属于（\quad）
\choicesline{A. 人机共生思想}{B. 新信息世界观}{C. 狭义工具论}{D. 计算思维}

\q 我国信息化发展的纲领性文件是（\quad）
\choices{A.《2006—2020 年国家信息化发展战略》}{B.《走近信息社会：中国信息社会发展报告 2010》}
{C.《第三次浪潮》}{D.《控制论》}

\qtypename{二、多项选择题}{3 题，每题 4 分，共 12 分；每题有两个或两个以上正确项}
\q 信息化的“四化”包括（\quad）
\choicesline{A. 智能化}{B. 电子化}{C. 全球化}{D. 群体化}

\q 信息化的“四性”包括（\quad）
\choicesline{A. 综合性}{B. 竞争性}{C. 渗透性}{D. 封闭性}

\q 信息社会的基本特征有（\quad）
\choicesline{A. 知识型经济}{B. 网络化社会}{C. 数字化生活}{D. 服务性政府}

\qtypename{三、判断题}{6 题，每题 2 分，共 12 分；正确的画“√”，错误的画“×”并写出改正后的表述}
\q 我国《2006—2020 年国家信息化发展战略》提出，到 2020 年综合信息基础设施基本普及、信息技术自主创新能力显著增强、国家信息安全保障水平大幅提高。\quad（\quad）\quad 改正：\fixblank

\q 信息社会就是互联网普及的社会，只要网络覆盖率达到一定水平就建成了信息社会。\quad（\quad）\quad 改正：\fixblank

\q 信息社会的战略资源是资本和能源。\quad（\quad）\quad 改正：\fixblank

\q 第四次信息技术革命是电子计算机和现代通信技术在信息工作中的应用。\quad（\quad）\quad 改正：\fixblank

\q 信息化的实现要以信息的数字化为必备条件。\quad（\quad）\quad 改正：\fixblank

\q 计算思维就是编写计算机程序的思维。\quad（\quad）\quad 改正：\fixblank

\qtypename{四、填空题}{6 空，每空 1 分，共 6 分}
\q 信息化特性可归纳为“四化”和“四性”，其中“四化”指智能化、电子化、全球化和\fillblanks[2.0cm]。

\q 信息社会以信息技术为基础，以\fillblanks[2.2cm]为支柱，以信息价值的生产为中心，以信息产品为标志。

\q 人类社会的第二次信息技术革命的标志是\fillblanks[2.6cm]的出现。

\q 国家信息中心提出的信息社会发展指数的英文缩写是\fillblanks[1.4cm]，其值在 0.3 以下属于\fillblanks[1.4cm]期。

\q 从狭义工具论转到\fillblanks[2.0cm]，是信息科学技术认识转变的重要内容。

\newpage
\qtypename{五、简答题}{4 题，每题 6 分，共 24 分}
\q 什么是信息化？简述信息化的“四化”与“四性”。\anslines{2.4cm}

\q 什么是信息社会？简述信息社会的基本特征。\anslines{2.4cm}

\q 简述国家信息中心对信息社会发展阶段的划分，并说明我国信息社会发展的不平衡状况。\anslines{2.4cm}

\q 简述科技界对信息技术认识的主要转变（4 个方面），并指出信息科学技术（含计算机技术）当前的重点发展方向。\anslines{2.4cm}

\newpage
\qtypename{六、综合题（论述）}{1 题，共 16 分}
\q 某县依托农业信息服务平台，把气象、土壤、市场行情等数据汇集起来向农户和合作社共享，同时推广测土配方施肥、病虫害在线诊断和农产品电商。请结合这一案例，分析信息技术如何改变农业生产与政府治理；要求用上信息产业、信息资源共享、数字鸿沟等概念。\anslines{5cm}

\end{document}
